\section{Analysis: Discriminability $\neq$ Reward Utility}\label{sec:discussion}

\paragraph{Unified root cause.}
The failures in Sections~\ref{sec:cosine_failure} and~\ref{sec:probe_trap} trace to a single geometric limitation.
$\dsyc$ encodes the \emph{direction} of belief change but not its \emph{evidential warrant}: valid revision and sycophantic capitulation both shift activations toward the user's position along the same direction.
The two reward families expose complementary faces of this blindness.
Cosine $\rone$ penalizes all belief change indiscriminately; the resulting mode collapse (F1c, F1d) reflects excessive punishment of any activation shift along $\dsyc$.
The binary probe $\rone$ rewards all resistance regardless of input validity; the stubbornness trap ($\tof = 0$ by step 500) reflects excessive permissiveness toward activation stasis.
These are not independent design flaws but dual consequences of a direction that conflates the target behavior (resisting invalid challenges) with its complement (refusing valid corrections).
Cross-model evidence confirms the limitation extends beyond Qwen3-8B: Llama-3-8B independently achieves AUROC $0.779 \pm 0.019$ on the same contrastive pairs with the same monotonic evidence gradient, so discriminability generalizes across architectures while the semantic gap between discriminability and reward utility persists.
This pattern contrasts with reasoning faithfulness and safety directions~\citep{zhang2026latentgrpo}, where the behavior to suppress is geometrically distinct from the behavior to preserve along the relevant direction.
For sycophancy, the target and its complement occupy the same activation region; they differ only in whether external evidence warrants the belief change.

\paragraph{GRPO stubbornness across all formulations.}

\begin{table}[t]
\centering
\caption{Comparison of all sycophancy calibration methods on Qwen3-8B. $\sycon$ is the harmonic mean of $\tof$ and $\nof$; higher indicates better calibration. Best result in \textbf{bold}. $^\dagger$Peak at step 4500; regresses to 0.114 by step 5000.}
\label{tab:main_comparison}
\begin{tabular}{llcccc}
\toprule
Method & Type & $\sycon$ $\uparrow$ & $\tof$ $\uparrow$ & $\nof$ $\uparrow$ & Step \\
\midrule
Vanilla & -- & 0.134 & 0.072 & 0.945 & -- \\
$\rthree$-only GRPO & RL & 0.134$^\dagger$ & 0.072 & 0.961 & 4500 \\
Cosine $\rone$ (F1d) & RL & collapse & -- & -- & -- \\
Binary Probe $\rone$ & RL & 0.000 & 0.000 & 0.995 & 500 \\
ActAdd ($\alpha{=}5$) & Inference & 0.054 & 0.028 & 0.961 & -- \\
\textbf{ACT} & SFT & \textbf{0.301} & \textbf{0.182} & 0.867 & 8000 \\
\bottomrule
\end{tabular}
\end{table}

Table~\ref{tab:main_comparison} summarizes all methods evaluated on Qwen3-8B.
$\rthree$-only GRPO removes $\dsyc$ entirely and relies on the NLI-based behavioral reward alone.
If the failures in Sections~\ref{sec:cosine_failure} and~\ref{sec:probe_trap} stemmed solely from the representation reward, this ablation should recover calibration quality.
It does not.
The $\sycon$ trajectory shows partial recovery that peaks at vanilla: 0.074 at checkpoint 1000, 0.085 at checkpoint 2000, 0.095 at checkpoint 3000, 0.114 at checkpoint 4000, a transient spike to 0.134 at checkpoint 4500---approximately matching vanilla in $\sycon$---followed by regression to 0.114 at checkpoints 5000 and 5500.
The best performance $\rthree$-only GRPO achieves is recovery to the untrained vanilla level; behavioral reward alone cannot push the model beyond its pre-training baseline.
ActAdd~\citep{turner2023actadd}, the inference-time baseline, achieves only $\sycon = 0.054$, confirming that activation-level intervention without weight updates cannot resolve the calibration problem.
$\rthree$-only GRPO exhibits the same stubbornness pattern as the $\dsyc$ approaches (high $\nof$, near-zero $\tof$) at milder severity.
The sparse scalar reward per completion cannot teach the model \emph{which} belief changes are warranted: it collapses an entire multi-turn exchange into one number that does not distinguish legitimate resistance from stubborn refusal.

\paragraph{Why SFT succeeds where GRPO fails.}
ACT~\citep{consistency2025bctact} achieves $\sycon = 0.301$ at checkpoint 8000, more than $2{\times}$ the vanilla baseline and above every RL and inference-time variant (Table~\ref{tab:main_comparison}).
The ACT trajectory is not monotonic: $\sycon = 0.255$ at checkpoint 2000, dips to 0.214 at checkpoint 4000 ($-4.1$ percentage points as $\tof$ drops and $\nof$ rises), recovers to 0.262 at checkpoint 6000, and reaches 0.301 at checkpoint 8000.
ACT also experiences stubbornness episodes, yet it recovers from them.
The critical difference is supervision density.
SFT provides a corrective gradient at every token position in every training example: when the model begins to resist valid corrections, per-token cross-entropy loss immediately penalizes the incorrect tokens and supplies gradient information about \emph{which} outputs need to change.
GRPO provides one scalar per completion; once the policy locks into stubbornness, that scalar reinforces the behavior because consistent non-flipping collects high reward from both $\rone$ and $\rthree$.
Dense per-token loss can break out of stubbornness equilibria that sparse per-completion reward cannot escape.
This recovery capacity is the practical signature of the dense-versus-sparse supervision gap.
ACT is also far more efficient: its checkpoint 500 ($\sycon = 0.170$) already exceeds the plateau that $\rthree$-only GRPO reaches after 5000 steps of training.

\paragraph{Future directions.}
These results point to several open problems.
First, condition-aware reward design could jointly encode belief-change direction and evidential warrant by conditioning on whether the input challenge provides valid grounds for revision.
Second, hybrid pipelines that apply SFT to establish the valid/invalid distinction as a warm start, followed by RL for broader distributional coverage, may combine the calibration quality of dense supervision with the exploration benefits of reinforcement learning.
Third, process reward models that provide per-step feedback within the RL framework could bridge the supervision density gap and retain online optimization.
Beyond sycophancy, the principle that discriminability does not guarantee reward utility applies to any behavioral attribute where the target and its complement share the same geometric direction~\citep{skalse2022hacking}.
Whether this principle extends to other alignment-relevant behaviors, such as hallucination and instruction-following, is an important open question.